\chapter{同调镜像对称初探}
\SourcePage{95}
这里的内容比简单地说“导出范畴就是 Fukaya 范畴”更微妙。以下始终设 $(X,\omega)$ 为紧 Calabi–Yau 流形。
\section{概览}
此前讨论的是闭弦，图像如图~\ref{fig:6.1} 所示。拓扑扭曲产生 A 模型，即 Gromov–Witten 理论（GW），以及 B 模型，即 Hodge 结构变分（VHS）。同调镜像对称则是开弦的版本，其图像如图~\ref{fig:6.2} 所示。
\OriginalFigure{fig-6-1.pdf}{0.2946}{一条闭弦随时间传播，并分裂成两条弦。}{fig:6.1}
\OriginalFigure{fig-6-2.pdf}{0.4200}{一条开弦随时间传播，边界条件为 $L_0,L_1$。}{fig:6.2}
\SourcePage{96}
\OriginalFigure{fig-6-3.pdf}{0.3746}{圆盘到流形的映射 $u:\Delta\to X$，其边界各段映到 Lagrange 子流形 $L_i$；可把它看作一条开弦的产生与消亡。}{fig:6.3}
\OriginalFigure{fig-6-4.pdf}{0.3081}{映射 $u:\Delta\to X$ 有三个边界标记点，各段边界分别映到三个 Lagrange 子流形。计数这样的映射，将在子流形的交点集合上给出代数结构。}{fig:6.4}
于是，我们有一些满足边界条件 $L_1,L_2$ 的开弦在空间中传播。扭曲后得到 A 模型，在亏格零时由 Fukaya 范畴 $\mathcal F(X)$ 给出。\footnote{还没有人敢于在高亏格情形中作出这一定义。}B 模型由导出范畴 $D^b(\OO_X)$ 给出。
\subsection{A 模型与 B 模型}
从图像上看，弦像图~\ref{fig:6.3} 那样产生和消亡。这给出交点之间的某种量子修正。一般情形则是图~\ref{fig:6.4} 那样的多边形。

边界条件给出“对象”的集合 $\Obj(\mathcal A)$，边界条件的交点给出“态射”的集合 $\Hom(\mathcal A)$。弦的传播由态射的复合运算 $m_k$ 给出，图~\ref{fig:6.3}、\ref{fig:6.4} 为这些 $m_k$ 赋予代数结构。

设 $(X,\omega)$ 为复维数 $d$ 的辛流形。当 $\mathcal A=\mathcal F(X)$ 时，对象是\Term{lagrangeziliuxing}{Lagrange 子流形}，即满足 $\dim_{\RR} L=d$、$\omega|_L=0$ 的子流形 $L\subset X$。传播中的弦对应于映射 $u:\Delta\to X$，其中 $\Delta=\{|z|\le1\}\subset\CC$ 为闭单位圆盘。若 $L,L'$ 横截相交，则
\Eq{6.1}{\operatorname{hom}(L,L')=\CC^{L\cap L'}.}
\SourcePage{97}
\OriginalFigure{fig-6-5.pdf}{0.3443}{定义 $m_k$ 时所计数的圆盘类型。}{fig:6.5}
实际上，还希望在这些 Lagrange 子流形上配备 $U(1)$ 局部系统；交点处的系数不再只是 $\CC$，而来自对应纤维之间的同态。运算
\Eq{6.2}{m_k:\Hom(L_{k-1},L_k)\otimes\cdots\otimes\Hom(L_0,L_1)\longrightarrow\Hom(L_0,L_k)}
由计数图~\ref{fig:6.5} 中的映射 $u:\Delta\to X$ 给出，且 $u(\partial\Delta)\subset L_0\cup\cdots\cup L_k$。

B 模型中使用范畴 $D^b(\OO_X)$，即凝聚层的有界导出范畴。其对象由复形
\Eq{6.3}{F^\bullet=(\cdots\longrightarrow F^{-r}\longrightarrow F^{-r+1}\longrightarrow\cdots)}
表示，要求上同调有界：\footnote{上标 $b$ 就表示有界。}
\Eq{6.4}{h^i(F^\bullet)=0\qquad\text{当 }|i|\gg0.}
这为什么重要？闭弦的故事主要是计数几何，关注它的人相对有限；开弦的故事则带来许多人关心的新对象。Lagrange 子流形是辛几何的重要目标，而导出范畴在几何表示论、代数几何等领域都非常重要。
\subsection{代数结构}
定义域的模空间，即稳定圆盘的模空间，具有递归结构。
\examnum{1}\begin{example}\label{ex:6.1}
考虑边界上有四个标记点的圆盘。计数其上的映射给出 $m_3$。它能以多种方式退化，其中两种见图~\ref{fig:6.6}。把标记点 $0$ 看作输出，其余点看作输入。这些退化可用带根的度量带状树编码，\footnote{带状图是指每个顶点处的边带有循环次序的图。这里的带状树总可嵌入平面。}对应的树也画在图~\ref{fig:6.6} 中。
\end{example}
对于两两横截的 $L_0,\ldots,L_k$，可通过对圆盘到 $X$ 的 $J$ 全纯映射作虚拟计数，定义
\Eq{6.5}{m_k:\Hom(L_{k-1},L_k)\otimes\cdots\otimes\Hom(L_0,L_1)\longrightarrow\Hom(L_0,L_k)[2-k].}
此模空间的虚拟维数为
\Eq{6.6}{(k-2)+d_0-\sum_{i=1}^kd_i.}
\SourcePage{98}
\OriginalFigure{fig-6-6.pdf}{0.8056}{四标记点圆盘的两种退化，以及相应的带根带状树。}{fig:6.6}
在适当条件下，交点带有 $\ZZ$ 分次，这就是移位 $[2-k]$ 的含义。研究高维模空间的边界，便得到 $m_k$ 之间的代数关系。
\examnum{2}\begin{example}\label{ex:6.2}
\begin{enumerate}[label=(\alph*)]
\item 第一个关系为 $m_1^2=0$，即 $m_1$ 是微分。取两个 Lagrange 子流形 $L_0,L_1$；令交点次数差为 $2$，便可看到条带断裂的边界，从而说明 $m_1$ 是
\Eq{6.7}{\bigoplus_{L,L',d}\Hom^d(L,L')}
上的微分。
\item 有关系
\Eq{6.8}{m_1m_2=m_2(m_1\otimes\id+\id\otimes m_1),}
它表达了 $m_1$ 相对于 $m_2$ 的分次 Leibniz 法则。
\item $m_2$ 本身未必结合，$m_3$ 则提供结合性的同伦。具体地，
\Eq{6.9}{\begin{aligned}m_2(\id\otimes m_2-m_2\otimes\id)={}&m_1m_3\\&+m_3(m_1\otimes\id\otimes\id+\id\otimes m_1\otimes\id+\id\otimes\id\otimes m_1).\end{aligned}}
右侧第一项是上边缘，第二项在 $m_1$-上闭链上为零。
\item 一般公式为
\Eq{6.10}{\sum_{\substack{p\ge1,\ q\ge0\\p+q\le k}}(-1)^{|a_1|+\cdots+|a_q|-q}\,m_{k-p+1}(a_k,\ldots,a_{p+q+1},m_p(a_{p+q},\ldots,a_{q+1}),a_q,\ldots,a_1)=0.}
\end{enumerate}
这些称为 $A_\infty$ 关系。
\end{example}
\SourcePage{99}
\subsection{Fukaya 范畴}
假定 $c_1(X)=0$，$X$ 紧，所有 $L_i$ 均定向、带有自旋结构与分次，并且无障碍。\footnote{例如，不存在以 $L_i$ 为边界的非恒值全纯圆盘时，这一条件成立。}此时可对两两横截的 Lagrange 子流形 $L_0,\ldots,L_k$ 定义 $m_k$。自然的系数环是 Novikov 环：
\Eq{6.11}{\Lambda_{\mathrm{nov}}=\left\{\sum_{i\ge0}a_iq^{\lambda_i}:\lambda_i\in\RR,\ \lambda_i\to+\infty,\ a_i\in\ZZ\right\}.}
也可以取 $a_i\in\QQ$ 或 $\CC$，此时得到\Term{novikovyu}{Novikov 域}。$a_i$ 来自全纯圆盘的计数，而 $\lambda_i$ 是相应圆盘的辛面积：
\Eq{6.12}{\int_\Delta u^*\omega.}
先验而言，这只给出 $A_\infty$ 预范畴，因为 $\Hom(L,L)$ 在横截交点的定义下没有意义。可用 Hamilton 扰动解决这一问题：对 $L,L'$，选取 $H_t:X\to\RR$，$t\in[0,1]$，其 Hamilton 流 $\varphi_t$ 满足 $\varphi_1(L')\pitchfork L$。可以证明，所得 $A_\infty$ 范畴 $\mathcal F(X)$ 在典范的 $A_\infty$ 拟等价意义下良定义。现在正式给出定义。
\thmnum{definition}{1}\begin{definition}\label{def:6.1}
一个非幺 $A_\infty$ 代数，是 $\ZZ$ 分次 $K$ 向量空间
\Eq{6.13}{A=\bigoplus_{p\in\ZZ}A^p,}
连同对每个 $k\ge1$ 的次数零分次 $K$ 线性映射
\Eq{6.14}{m_k:A^{\otimes k}\longrightarrow A[2-k],}
它们满足 $A_\infty$ 关系。
\end{definition}
\thmnum{definition}{2}\begin{definition}\label{def:6.2}
一个非幺 $A_\infty$ 范畴 $\mathcal A$，由对象集合 $\Obj(\mathcal A)$，对每对对象 $X_0,X_1$ 的 $\ZZ$ 分次 $K$ 向量空间
\Eq{6.15}{\Hom(X_0,X_1),}
以及对每个 $k\ge1$ 的映射
\Eq{6.16}{m_k:\Hom(X_{k-1},X_k)\otimes\cdots\otimes\Hom(X_0,X_1)\longrightarrow\Hom(X_0,X_k)[2-k]}
组成，且满足 $A_\infty$ 关系。
\end{definition}
对 $A_\infty$ 范畴取其上同调范畴 $H(\mathcal F(X))$，就得到通常的范畴：对象不变，态射空间取关于 $m_1$ 的上同调；此时 $m_1=0$。
\begin{custom}{事实 9}\hypertarget{fact:9}{}$H(\mathcal F(X))$ 是幺范畴，即对每个 $L$，都存在某个 $e_L\in\Hom^0(L,L)$，满足 $m_1(e_L)=0$，它对 $m_2$ 起单位元作用，而以它为输入时，所有 $k>2$ 的 $m_k$ 均为零。\end{custom}
这样的范畴称为上同调幺范畴，或 $c$ 幺范畴。它稍有不便，不过有一种技巧：$c$ 幺范畴 $\mathcal F(X)$ 可以证明与某个严格幺 $A_\infty$ 范畴 $A_\infty$-拟等价。
\SourcePage{100}
\subsection{“导出”Fukaya 范畴}
这里也许会产生疑问：如何把这个 $A_\infty$ 范畴与一个三角导出范畴对应？Kontsevich 的办法是人为加入复形，把 $\mathcal A=\mathcal F(X)$ 扩张为一个“三角 $A_\infty$ 范畴” $\operatorname{Tw}\mathcal F(X)$。
具体地，对 $L_1,\ldots,L_k$，取 $\Delta=(\delta_{ij})$，$\delta_{ij}\in\Hom^1(L_j,L_i)$，满足：
\begin{itemize}
\item 严格下三角：当 $j\ge i$ 时 $\delta_{ij}=0$；
\item Maurer–Cartan 方程：
\Eq{6.17}{\sum_{r\ge1}m_r(\Delta,\ldots,\Delta)=\sum_{\substack{i_0<\cdots<i_r\\r\ge1}}m_r(\delta_{i_r,i_{r-1}},\ldots,\delta_{i_1,i_0})=0.}
\end{itemize}
\Lecture{27}{2019年12月5日}
\subsection{同调镜像对称}
现在给出同调镜像对称的精确陈述。回忆，Fukaya 范畴的系数属于 Novikov 域；采用复系数，令 $K=\Lambda_{\mathrm{nov},\CC}$。紧 Calabi–Yau 流形 $(X,\omega)$ 给出镜像的极大退化族 $Y/S$，此处在 $S=\Spec K$ 上工作。
\begin{custom}{猜想 1}\hypertarget{conj:1}{}
在 $K$ 上，存在典范的 $A_\infty$ 范畴等价
\Eq{6.18}{\psi:D^b\mathcal F(X)=H^0(\operatorname{Tw}^{\Pi}\mathcal F(X))\xrightarrow{\sim}D^b(Y).}
\end{custom}
$\operatorname{Tw}^{\Pi}\mathcal F(X)$ 称为分裂闭包；换言之，它是在扭曲 Fukaya 范畴中加入幂等元的像所得到的完备化。
\section{导出范畴的微分分次增强与同调扰动引理}
\subsection{导出范畴的\texorpdfstring{Čech}{Čech}模型}
微分分次范畴，即 dg 范畴，是 $k>2$ 时 $m_k=0$ 的 $A_\infty$ 范畴。我们希望描述 $D^b(X)$ 的 $A_\infty$ 结构。关键在于，取上同调后仍可保留一个 $A_\infty$ 结构。\Term{tongdiaoraodongyinli}{同调扰动引理}（HPL）说明：即使从 dg 范畴出发，传到上同调后的结构一般也只是 $A_\infty$ 的。

假定 $X$ 光滑且射影。每个凝聚 $\OO_X$ 模都有有限的局部自由消解，故
\Eq{6.19}{D^b(X)=D^b(\operatorname{Coh}X)=\operatorname{Perf}(X).}
一般而言，拟凝聚层导出范畴中的完美复形就是紧对象；不过若从一开始就假定射影，所涉及的对象本来就都是紧的。

dg 增强如下：$D^b_{\mathrm{dg}}(X)$ 的对象是局部自由层的有界复形，态射才是需要说明的部分。记 $X$ 的有限仿射覆盖为
\Eq{6.20}{\mathscr U=\{U_i\}.}
令
\Eq{6.21}{\Hom(E^\bullet,F^\bullet)=\bigoplus_n\Hom^n_{D^b_{\mathrm{dg}}(X)}(E^\bullet,F^\bullet),}
其中
\Eq{6.22}{\Hom^n_{D^b_{\mathrm{dg}}(X)}(E^\bullet,F^\bullet)=\bigoplus_{p+q=n}\check C^p(\mathscr U,\mathcal Hom^q_{\OO_X}(E^\bullet,F^\bullet)).}
这里记
\Eq{6.23}{\mathcal Hom_{\OO_X}(E^\bullet,F^\bullet[q])=\mathcal Hom^q_{\OO_X}(E^\bullet,F^\bullet),}
\SourcePage{101}
$\check C^p$ 为 Čech $p$ 上链：
\Eq{6.24}{\check C^p(\mathscr U,\mathcal G)=\bigoplus_{i_0<\cdots<i_p}\Gamma(U_{i_0}\cap\cdots\cap U_{i_p},\mathcal G).}
这构成 dg 范畴，微分 $m_1$ 由 Čech 微分和内部 Hom 微分给出。
对于
\Eq{6.25}{f=(f_{n,i_0,\ldots,i_p})\in\check C^p(\mathscr U,\mathcal Hom^q(E^\bullet,F^\bullet)),}
\Eq{6.26}{g=(g_{n',i_0,\ldots,i_{p'}})\in\check C^{p'}(\mathscr U,\mathcal Hom^{q'}(F^\bullet,G^\bullet)),}
复合为
\Eq{6.27}{m_2(g,f)=(-1)^{\deg f}g\circ f,}
其中
\Eq{6.28}{g\circ f=\left((-1)^{q'p}g_{n+q,i_0,\ldots,i_{p'}}\circ f_{n,i_{p'},\ldots,i_{p+p'}}\right)\big|_{U_{i_0}\cap\cdots\cap U_{i_{p+p'}}}.}
\subsection{同调扰动引理}
\thmnum{theorem}{1}\begin{theorem}\label{thm:6.1}
对任意定义在域上的 $A_\infty$ 范畴 $\mathcal A$，其上同调 $H(\mathcal A)$ 上存在 $m_1=0$ 的 $A_\infty$ 范畴结构，称为极小 $A_\infty$ 范畴，并且有 $A_\infty$ 拟等价 $H(\mathcal A)\to\mathcal A$。
\end{theorem}
\begin{proof}
该构造依赖于对每对对象 $X,Y$ 选择投影算子 $\Pi$ 与同伦 $H$：
\Eq{6.29}{\Pi:\Hom(X,Y)\longrightarrow\Hom(X,Y),}
\Eq{6.30}{H:\Hom(X,Y)\longrightarrow\Hom(X,Y)[-1],}
其中 $\Pi$ 次数为零，$\Pi^2=\Pi$，并满足
\Eq{6.31}{\Pi-\id=m_1H+Hm_1.}
这将给出所需的拟等价。取 $\Pi$ 为到上同调类的闭代表元所组成子空间的投影。

得到一个新 $A_\infty$ 范畴 $\Pi\mathcal A$，其对象不变，态射为
\Eq{6.32}{\Hom_{\Pi\mathcal A}(X,Y)=\Pi(\Hom_{\mathcal A}(X,Y))\subset\Hom_{\mathcal A}(X,Y),}
故 $H(\Pi\mathcal A)=H(\mathcal A)$。定义运算如下：
\Eq{6.33}{m_1^\Pi=m_1,}
\Eq{6.34}{m_2^\Pi=\Pi\circ m_2,}
\Eq{6.35}{m_k^\Pi=\sum_{\substack{T\text{ 为带根带状树}\\T\text{ 有 }k\text{ 个叶子}}}m_{k,T}\qquad(k\ge3).}
这些规则可看作 Feynman 规则：每个有 $r+1$ 条邻接边的内部顶点放置 $m_r$，内部边放置 $H\circ$，根端放置 $\Pi\circ$。
\examnum{3}\begin{example}\label{ex:6.3}
取 $k=6$。图~\ref{fig:6.7} 给出求和中的一棵带根带状树，其贡献为
\Eq{6.36}{m_{6,T}(a_1,\ldots,a_6)=\Pi m_3(a_1,Hm_2(Hm_3(a_2,a_3,a_4),a_5),a_6).}
\end{example}
\SourcePage{102}
\OriginalFigure{fig-6-7.pdf}{0.5328}{对 $m_6^\Pi$ 的定义作出贡献的一棵带根带状树。}{fig:6.7}
\thmnum{lemma}{2}\begin{lemma}\label{lem:6.2}
$m_k^\Pi$ 满足 $A_\infty$ 关系。
\end{lemma}
\begin{Note}{6.1}存在 $A_\infty$ 函子
\Eq{6.37}{(i_k):\Pi\mathcal A\longrightarrow\mathcal A,}
它与 $H(\mathcal A)$ 上由 $m_k^\Pi$ 给出的 $A_\infty$ 结构相容。$i_1$ 为包含映射；其余 $i_k$ 用类似的树构造定义：
\Eq{6.38}{i_k=\sum_{\substack{T\text{ 为带根带状树}\\T\text{ 有 }k\text{ 个叶子}}}i_{k,T}\qquad(k\ge2),}
其中 $i_{k,T}$ 与 $m_{k,T}$ 的定义相同，只把根端的 $\Pi$ 换成 $H$。
\end{Note}
\thmnum{lemma}{3}\begin{lemma}\label{lem:6.3}
$(i_k)$ 是 $A_\infty$ 函子。
\end{lemma}
\end{proof}
\begin{proof}[引理~\ref{lem:6.2} 的证明概要]
考虑
\Eq{6.39}{\mu_k=\sum_{\substack{T\text{ 为带状树}\\e\text{ 为一条边}}}\mu_{k,(T,e)}.}
其中 $\mu_{k,(T,e)}$ 与 $m_{k,T}$ 的构造相同，但在内部边 $e$ 上，把 $H$ 换成 $\Pi-\id=m_1H+Hm_1$；对于外部边，则放置 $m_1$。有两种整理方式：
\begin{enumerate}[label=(\roman*)]
\item 使用 $\Pi-\id$，按边 $e$ 分类：
\Eq{6.40}{\mu_k=\underbrace{\mu_k^{m_1}}_{e\text{ 为外部边}}+\underbrace{\mu_k^\Pi-\mu_k^{\id}}_{e\text{ 为内部边}}.}
\item 使用 $m_1H+Hm_1$，按与 $e$ 相邻的内部顶点分类。
\end{enumerate}
再对各 $m_r$ 应用 $A_\infty$ 关系，得到新的树求和，其中每项出现在 $\mu_k^{\id}$ 中。核对分次符号后，$\mu_k=-\mu_k^{\id}$。由\eqref{eq:6.40}得
\Eq{6.41}{\mu_k^{m_1}+\mu_k^\Pi=0.}
这恰好是 $m_k^\Pi$ 的 $A_\infty$ 关系。
\end{proof}
\SourcePage{103}
\section{Polishchuk 的刚性定理}
\subsection{约化到代数\texorpdfstring{$A_L$}{AL}}
先设 $X$ 为光滑射影簇，$L\in D^b(X)$ 为丰线丛。在 $D^b_{\mathrm{dg}}(X)$ 中，取以 $L^n$（$n\in\ZZ$）为对象的全子范畴 $\langle L\rangle$。在 $D^b(X)=H(D^b_{\mathrm{dg}}(X))$ 上，$m_2$ 由下列代数上的杯积给出：
\Eq{6.42}{A_L=\bigoplus_{p,q\in\ZZ}H^p(X,L^q).}
注意，$A_L$ 包含齐次坐标环：
\Eq{6.43}{A_L\supseteq R_L=\bigoplus_{n\ge0}H^0(X,L^n).}
$A_L$ 带有极小\footnote{这表示 $m_1=0$。}$A_\infty$ 代数结构，而 $\langle L\rangle$ 分裂生成 $D^b(X)$。因此，$A_L$ 连同其内部幂次分次及 $A_\infty$ 代数结构，刻画了 $D^b(X)$ 的 $A_\infty$ 范畴结构。
\Lecture{28}{2019年12月10日}
\subsection{\texorpdfstring{$A_L$ 上的 $A_\infty$ 结构}{AL 上的 A∞ 结构}}
于是有如下问题。
\setcounter{prob}{0}\begingroup\renewcommand{\theprob}{\arabic{prob}}\begin{problem}\label{problem:6.1}
除 $m_2$ 之外，$A_L$ 上的 $A_\infty$ 结构还保存了多少信息？
\end{problem}\endgroup
\begin{Note}{6.2}
设 $X$ 为特征零域上的 $d$ 维光滑射影 Calabi--Yau 簇，且 $H^p(X,\OO_X)=0$ 对 $0<p<d$ 成立，则有
\Eq{6.44}{\begin{aligned}
A_L&=\bigoplus_{p\ge0,\ q\in\ZZ}H^p(X,L^{\otimes q})\\
&=\bigoplus_{q\ge0}H^0(X,L^{\otimes q})
\oplus\bigoplus_{q\le0}\underbrace{H^d(X,L^{\otimes q})}_{\simeq H^0(X,L^{\otimes(-q)})^*}.
\end{aligned}}
\end{Note}
\begin{Note}{6.3}线丛 $L$ 的幂次给出了一个内部次数分次。\end{Note}
\thmnum{theorem}{4}\begin{theorem}[Polishchuk]\label{thm:6.4}
设 $X$ 为 $d$ 维射影簇，$d\ge1$，$L$ 为甚丰线丛，并且对所有 $p\in\ZZ$ 及 $q\ne0,d$，都有 $H^q(X,L^{\otimes p})=0$。则：
\begin{enumerate}[label=(\roman*)]
\item 在严格 $A_\infty$ 同构（这里指 $(f_n)$ 满足 $f_1=\id$）与缩放的意义下，$A_L$ 上存在唯一的保持内部次数、扩张既定 $m_2$ 的非平凡极小 $A_\infty$ 结构。非平凡表示它不能通过这样的同构化为所有 $k>2$ 的 $m_k$ 均为零的结构。
\item 两个给定结构之间的严格 $A_\infty$ 同构，在同伦意义下唯一。
\end{enumerate}
\end{theorem}
\begin{Note}{6.4}
这说明，同调镜像对称的内容在某种意义下可以集中到一个很小的部分。刚性的作用在于：要证明同调镜像对称，只需把两边的环对应起来，再验证这少量额外的信息一致。
\end{Note}
\subsection{\texorpdfstring{$A_\infty$}{A∞}结构与 Hochschild 上闭链}
设 $A=\bigoplus_{k\in\ZZ}A^k$ 为分次代数，$(m_n)$ 与 $(m'_n)$ 是其两个 $A_\infty$ 扩张，其中 $m_1=m'_1=0$，$m_2=m'_2$ 为给定的乘法。如果对所有 $i<n$，都有 $m_i=m'_i$，那么：
\begin{enumerate}[label=(\roman*)]
\item $c=m_n-m'_n$ 是一个 Hochschild $n$ 上闭链，其上边缘算子为
\Eq{6.45}{\begin{aligned}
\delta c(a_1,\ldots,a_{n+1})={}&(-1)^{n\deg a_1}a_1c(a_2,\ldots,a_{n+1})\\
&+\sum_{j=1}^{n}(-1)^j c(a_1,\ldots,a_ja_{j+1},\ldots,a_{n+1})\\
&+(-1)^{n+1}c(a_1,\ldots,a_n)a_{n+1}.
\end{aligned}}
这是因为，把两个 $A_\infty$ 关系相减后，只剩下这些项。
\item 用保持较低阶数据的严格 $A_\infty$ 同构改变 $m_n$，只会给 $c$ 加上一个 Hochschild 上边缘。
\end{enumerate}
\SourcePage{104}
\subsection{这里的具体情形}
现在 $A=R\oplus M$，其中
\Eq{6.46}{R=\bigoplus_{q\ge0}H^0(X,L^q),\qquad\deg R=0,}
\Eq{6.47}{M=\bigoplus_{q\le0}H^d(X,L^q),\qquad\deg M=d,}
且
\Eq{6.48}{m_n:A^{\otimes n}\longrightarrow A[2-n].}
考察次数。如果输入 $a_i$ 中有 $k$ 个属于 $M$，则输入总次数为 $kd$，其中 $0\le k\le n$，输出次数为 $kd+2-n$。由于 $A$ 只有次数 $0$ 与 $d$，必有
\Eq{6.49}{kd+2-n=0\qquad\text{或}\qquad(k-1)d+2-n=0.}
因此，令 $l=k$ 或 $k-1$，则在非零的情形中 $n=ld+2$，$0\le l\le n$。两个结构第一次不同之处，$c=m_n-m'_n$ 是内部次数零、上同调次数 $2-n$ 的 Hochschild 上闭链，记作
\Eq{6.50}{c\in C^n_{0,n-2}=C^{ld+2}_{0,ld}.}
下标依次记录内部次数与上同调次数的相反数。
\subsection{标准 bar 消解、cobar 复形与 Hochschild 上同调}
设 $A$ 为结合 $k$ 代数，其中 $k$ 是交换幺环，并给定满的增广映射 $\epsilon:A\twoheadrightarrow k=A/A_+$。考虑复形
\Eq{6.51}{\cdots\longrightarrow A_+\otimes A_+\otimes A
\longrightarrow A_+\otimes A\longrightarrow A\xrightarrow{\epsilon}k\longrightarrow0.}
取 $a_1,\ldots,a_n\in A_+$，$a_{n+1}\in A$，其微分由相邻因子的乘法给出：
\Eq{6.52}{d(a_1\otimes\cdots\otimes a_{n+1})=
\sum_{i=1}^{n}(-1)^i a_1\otimes\cdots\otimes a_ia_{i+1}\otimes\cdots\otimes a_{n+1}.}
\begin{custom}{事实 10}\hypertarget{fact:10}{}这是一个复形，即 $d^2=0$。\end{custom}
当 $A$ 有单位元时，此增广复形正合。这称为 \Term{bargouzao}{bar 构造}。\footnote{Eilenberg 与 Mac Lane 最初写出这一构造时，用竖线 $|$ 简记张量积；名称 bar 正由此而来。}
对 $A$ 双模 $M$，令
\Eq{6.53}{C^n(A,M):=\Hom_k(A_+^{\otimes n},M).}
它与上边缘算子 $\delta$ 构成复形，称为 cobar 复形。其上同调为
\Eq{6.54}{HH^n(A,M):=H^n(C^\bullet(A,M),\delta),}
\Eq{6.55}{HH^n(A):=HH^n(A,A),}
称为 \Term{Hochschildshangtongdiao}{Hochschild 上同调}。内部次数与上同调次数给出两重分次，例如记为
\Eq{6.56}{HH^n_{0,n-2}(A).}
\begin{Note}{6.5}一般的 $A_\infty$ 代数也有 Hochschild 同调与上同调，不过这里不需要使用它们。\end{Note}
\SourcePage{105}
\subsection{刚性定理的证明}
\thmnum{proposition}{5}\begin{proposition}\label{prop:6.5}
对上述 $A=R\oplus M$，有：
\begin{enumerate}[label=(\roman*)]
\item 对所有 $l\ge1$ 及 $i<l(d+2)$，
\Eq{6.57}{HH^i_{0,ld}(A)=0.}
\item $\dim HH^{d+2}_{0,d}(A)\le1$。
\end{enumerate}
\end{proposition}
\thmnum{corollary}{6}\begin{corollary}\label{cor:6.6}
两个 $A_\infty$ 扩张 $(m_n)$ 与 $(m'_n)$ 严格 $A_\infty$ 同构，当且仅当
\Eq{6.58}{[m_{d+2}-m'_{d+2}]=0\in HH^{d+2}_{0,d}(A).}
\end{corollary}
\begin{proof}
逐阶比较并消去差异。在当前第一次不同的阶数，有 $[m'_n-m_n]\in HH^{ld+2}_{0,ld}(A)$，其中 $0\le l\le n$。当 $l>1$ 时，$ld+2<l(d+2)$，故由命题~\ref{prop:6.5}，这个类自动为零。对于 $n>2$，只需另行考察 $l=1$，也就是 $m'_{d+2}-m_{d+2}$。
\end{proof}
\begin{proof}[定理~\ref{thm:6.4} 的证明]
由命题~\ref{prop:6.5} 的第(ii)项与推论~\ref{cor:6.6}，或者所有高阶运算都能消去，或者非平凡结构在缩放意义下唯一。
\end{proof}
\begin{Note}{6.6}
对光滑 $X$ 的典范结构，$[m_{d+2}]\ne0$，所以出现的是第二种可能。其依据是考察一点的理想层并取 Koszul 消解，从而得到 $m_{d+2}\ne0$。
\end{Note}
\begingroup\setstretch{1.15}
\setlength{\abovedisplayskip}{1pt}
\setlength{\belowdisplayskip}{1pt}
\setlength{\abovedisplayshortskip}{0pt}
\setlength{\belowdisplayshortskip}{1pt}
\begin{proof}[命题~\ref{prop:6.5} 的证明]
由 $A=R\oplus M=k\oplus R_+\oplus M$，对固定的 $l$，可写成
\Eq{6.59}{C^i:=C^i_{0,ld}=\underbrace{C^i(0)}_{\text{输出次数 }0}
\oplus\underbrace{C^i(d)}_{\text{输出次数 }d}.}
这些空间分别包含于
\Eq{6.60}{C^i(0)\subset\Hom(A_+^{\otimes i},R),}
\Eq{6.61}{C^i(d)\subset\Hom(A_+^{\otimes i},M).}
上链的形状如下。对 $c\in C^i(0)$，
\Eq{6.62}{c:[T(R_+)\otimes M\otimes T(R_+)\otimes\cdots\otimes M\otimes T(R_+)]_i\longrightarrow R.}
这里 $T(R_+)$ 为定义在 $k$ 上的张量代数，下标 $i$ 表示属于 $H^*(L^*)$ 的因子总数。对于 $C^i(d)$，有类似的描述。
由此得到正合序列
\Eq{6.63}{0\longrightarrow C^\bullet(d)\longrightarrow C^\bullet\longrightarrow C^\bullet(0)\longrightarrow0.}
因此，只需证明：
\begin{enumerate}[label=(\roman*)]
\item 对 $i<l(d+2)$，$H^i(C^\bullet(0))=H^i(C^\bullet(d))=0$；
\item 当 $l=1$ 时，$H^{d+2}(C^\bullet(d))=0$ 且 $\dim H^{d+2}(C^\bullet(0))\le1$。
\end{enumerate}
先按输出在 $R$ 中的内部次数分解：
\Eq{6.64}{C^\bullet(0)=\prod_{j\ge0}C^\bullet(0)_j.}
微分 $\delta$ 的形式为
\Eq{6.65}{\delta(x_j)_{j\ge0}=\left(\sum_{j\le j'}\delta_{j,j'}x_j\right)_{j'\ge0}.}
\SourcePage{106}
相邻因子相乘时，$M$ 因子的数目不会增加。现在利用下式逐步归纳：
\Eq{6.66}{0\longrightarrow\operatorname{gr}_i C^\bullet(0)
\longrightarrow}
这意味着，可以在各个 $(C^\bullet(0)_j,\delta_{j,j})$ 的上同调上验证第(i)、(ii)项。它们可写作
\Eq{6.67}{(C^\bullet(0)_j,\delta_{j,j})=\Hom(K^\bullet_{l,j},R_j)[-l],}
其中 $K^\bullet_{l,j}$ 为下述复形的第 $j$ 个内部次数分量：
\Eq{6.68}{K^\bullet_l=B^\bullet(k,M,\ldots,M,k).}
这里使用某种推广的 bar 复形：
\Eq{6.69}{B^\bullet(M_1,\ldots,M_n)
=M_1\otimes T(R_+)\otimes M_2\otimes\cdots\otimes T(R_+)\otimes M_n.}
相应的 bar 复形计算可以对 $n$ 归纳。这就是此处使用的 \href{https://arxiv.org/abs/math/0302178}{Polishchuk 的命题 3.7}。把 $B^\bullet(M_1,\ldots,M_n)$ 看成一个双复形的总复形，其双次数为
\Eq{6.70}{\deg_{0/1}(x_1\otimes t_1\otimes\cdots\otimes t_{n-1}\otimes x_n)
=\sum_{\substack{i\equiv0/1\pmod2}}\deg(t_i),}
于是谱序列的 $E_1$ 页形如
\Eq{6.71}{\begin{aligned}
H^*(M_1\otimes T(R_+)\otimes M_2)\otimes T(R_+)
\otimes H^*(M_3\otimes T(R_+)\otimes M_4)\otimes\cdots.
\end{aligned}}
接下来使用以下引理。
\thmnum{lemma}{7}\begin{lemma}\label{lem:6.7}
若 $M_1$ 与 $M_2$ 在左分次 $R$ 模和右分次 $R$ 模的意义下均同构于 $M$，则
\Eq{6.72}{H^i(B^\bullet(M_1,M_2))=
\begin{cases}M,&i=-d-1,\\0,&\text{其他情形},\end{cases}}
其中第一个同构分别是在左、右分次 $R$ 模意义下的同构。
\end{lemma}
因此，$E_1$ 在适当移位后化为 $B^\bullet(M'_1,\ldots,M'_{n'})$，其中 $n'<n$。
\end{proof}\endgroup
